\documentclass[../cours_tle.tex]{subfiles}
\begin{document}
\setcounter{chapter}{8}
\chapter{Comprendre la limite d’une suite}
\label{ch:9}
\begin{objectif}
Spécialité. Passer de l'intuition à une définition exploitable. Prérequis : chapitres 1 et 2.
\end{objectif}
\section{Se rapprocher durablement}
\begin{definition}
$u_n\to\ell\in\R$ signifie que tout intervalle ouvert contenant $\ell$ contient tous les $u_n$ à partir d'un certain rang. En notation précise :
$\forall\varepsilon>0,\ \exists N\in\N,\ \forall n\ge N,\ |u_n-\ell|<\varepsilon$.
\end{definition}
Le rang $N$ dépend de la précision demandée, mais pas de $n$ ensuite. Une suite n'a pas à atteindre sa limite : $1/(n+1)$ ne s'annule jamais, mais tend vers $0$ car $n+1>1/\varepsilon$ suffit.
Une suite convergente peut osciller : $(-1)^n/(n+1)\to0$.
\begin{theoreme}[Unicité et bornitude]
Une suite convergente possède une unique limite finie et est bornée.
\end{theoreme}
\begin{demonstration}
Si $\ell\ne m$ étaient deux limites, choisir $\varepsilon=|\ell-m|/3$. Pour un rang assez grand, l'inégalité triangulaire donnerait $|\ell-m|\le|\ell-u_n|+|u_n-m|<2|\ell-m|/3$, contradiction.
Avec $\varepsilon=1$, les termes à partir de $N$ vérifient $|u_n|\le|\ell|+1$. Les termes antérieurs sont en nombre fini et ont un maximum de leurs valeurs absolues ; prendre le plus grand des deux majorants.
\end{demonstration}
La réciproque est fausse : $(-1)^n$ est bornée mais ne converge pas, car les termes pairs valent $1$ et les impairs $-1$.
\section{Limites infinies}
$u_n\to+\infty$ signifie : pour tout réel $A$, il existe $N$ tel que $n\ge N\Rightarrow u_n>A$. Pour $-\infty$, remplacer par $u_n<A$.
Une suite non majorée peut revenir sans cesse à zéro ; cela ne suffit donc pas pour tendre vers $+\infty$.
\begin{theoreme}
Toute suite croissante non majorée tend vers $+\infty$.
\end{theoreme}
\begin{demonstration}
Fixer $A$. La non-majoration fournit un indice $N$ avec $u_N>A$. La croissance assure ensuite $u_n\ge u_N>A$ pour tous les $n\ge N$.
\end{demonstration}
Par exemple, $u_{2n}=2n$, $u_{2n+1}=0$ est non majorée mais ne tend pas vers $+\infty$.
\section{Comparer pour conclure}
\begin{theoreme}[Encadrement]
Si, à partir d'un rang, $a_n\le u_n\le b_n$ et $a_n,b_n\to\ell$, alors $u_n\to\ell$.
Si $u_n\ge a_n$ et $a_n\to+\infty$, alors $u_n\to+\infty$.
\end{theoreme}
\begin{demonstration}
Pour le premier résultat, rendre simultanément $a_n>\ell-\varepsilon$ et $b_n<\ell+\varepsilon$. L'encadrement impose les mêmes bornes à $u_n$.
Pour le second, dès que $a_n>A$, on a aussi $u_n>A$.
\end{demonstration}
Si $a_n\le b_n$ et les deux convergent, alors $\lim a_n\le\lim b_n$. Sinon des voisinages disjoints des deux limites inverseraient l'ordre pour les grands indices. Une inégalité stricte peut devenir une égalité à la limite : $1/n>0$, mais sa limite est $0$.
\section{Une méthode et ses limites}
Pour $u_n=(2n+\sin n)/(n+1)$, écrire
$u_n-2=(\sin n-2)/(n+1)$, donc $|u_n-2|\le3/(n+1)\to0$.
La preuve fournit aussi une précision : $n+1>3/\varepsilon$ suffit.
Une courbe ou une table de valeurs suggère le résultat, mais ne remplace pas cette borne.
\begin{finchapitre}
\section{Utilité et suite du parcours}
Les définitions de convergence seront réutilisées pour les séries, les fonctions, les méthodes numériques et les probabilités. Le quantificateur « à partir d'un rang » sépare les comportements locaux des comportements asymptotiques.
\begin{retenir}
Convergente implique bornée, pas l'inverse. Non majorée n'implique pas limite infinie sans hypothèse supplémentaire.
\end{retenir}
\end{finchapitre}
\end{document}
